\chapter{Late Tang: Li Shangyin and Ambiguity}

\section{The late Tang mood}

The ninth century was a time of decline: factional struggle at court (the ``Niu--Li''
factions), eunuch control of emperors, rebellions in the provinces, ending with the Huang
Chao rebellion (875--884) and the fall of the dynasty in 907. Late Tang poetry is often
described as autumnal: nostalgic, sensuous, introspective, fond of history, ruins, and the
private world of love.

\begin{center}
\begin{tabular}{llp{0.5\textwidth}}
\toprule
Poet & Dates & Known for \\
\midrule
Du Mu \zht{杜牧} & 803--852 & Elegant quatrains on history and pleasure quarters;
  Yangzhou; nostalgic ``meditations on the past'' \\
Li Shangyin \zht{李商隱} & 813?--858 & Dense, allusive, ambiguous regulated octaves;
  ``Untitled'' love poems \\
Wen Tingyun \zht{溫庭筠} & 812?--870 & Ornate poems and early song lyrics (\emph{ci});
  bridge to the Song \\
\bottomrule
\end{tabular}
\end{center}

\section{Du Mu: history in a quatrain}

``Mooring on the Qinhuai'' \zht{泊秦淮}. The Qinhuai River runs through Nanjing, the
capital of the short-lived Southern Dynasties, famous for its entertainment quarter.

\begin{center}
\begin{tabular}{p{0.36\textwidth}p{0.54\textwidth}}
\toprule
\zht{煙籠寒水月籠沙，} & Mist veils the cold water, moonlight veils the sand; \\
\zht{夜泊秦淮近酒家。} & at night I moor on the Qinhuai, near the wine houses. \\
\zht{商女不知亡國恨，} & The singing girls don't know the grief of a fallen kingdom: \\
\zht{隔江猶唱後庭花。} & across the river they are still singing ``Flowers of the Rear Court''. \\
\bottomrule
\end{tabular}
\end{center}

\textbf{Analysis.} ``Flowers of the Rear Court'' was composed by the last ruler of the
Chen dynasty, who partied while the Sui armies came to destroy him (589). The girls sing
it innocently as a pop song; the poet hears it as a warning to the Tang. Line 1 has
internal parallelism (\zht{煙籠寒水} / \zht{月籠沙}): the veiling mist and moonlight are
the haze of oblivion. This is the subgenre of \zht{詠史} ``poems on history''.

\section{Li Shangyin \zht{李商隱}}

\begin{itemize}
  \item A gifted young man patronised by Linghu Chu, a leader of the Niu faction; after
        passing the \emph{jinshi} (837) he married the daughter of Wang Maoyuan, associated
        with the rival Li faction. The Niu faction regarded him as a turncoat; his
        official career stalled in minor provincial posts.
  \item His wife died around 851; many poems may mourn her.
  \item Known for \textbf{dense allusion}, \textbf{sensuous imagery} (silk, candles,
        jade, incense, mist) and \textbf{ambiguity}: about twenty poems are simply titled
        \zht{無題} ``Untitled'', and it is often impossible to tell whether a poem is
        about a lover, a Daoist priestess, his wife, a patron, or his career.
\end{itemize}

\begin{important}[Three ways readers have read Li Shangyin]
\begin{enumerate}
  \item \textbf{Biographical}: each poem conceals a specific love affair (with a palace
        lady, a Daoist nun), reconstructed by detective work.
  \item \textbf{Political allegory} (the \emph{Li Sao} tradition): the beloved is a patron or
        the emperor; longing = desire for office.
  \item \textbf{Open / symbolist}: the poems create a mood of longing and loss whose object
        is deliberately left indeterminate. Modern readers (and the Chinese modernists of
        the 20th century) favour this.
\end{enumerate}
\end{important}

\section{``The Ornamented Zither'' \zht{錦瑟}}

Probably the most debated poem in Chinese. The title is just the first two characters.

\begin{center}
\begin{tabular}{p{0.36\textwidth}p{0.54\textwidth}}
\toprule
\zht{錦瑟無端五十弦，} & The ornamented zither, for no reason, has fifty strings; \\
\zht{一弦一柱思華年。} & each string, each bridge, recalls the flowering years. \\
\zht{莊生曉夢迷蝴蝶，} & Master Zhuang, dreaming at dawn, lost himself in a butterfly; \\
\zht{望帝春心託杜鵑。} & King Wang entrusted his springtime heart to the cuckoo. \\
\zht{滄海月明珠有淚，} & On the blue sea, under a bright moon, pearls have tears; \\
\zht{藍田日暖玉生煙。} & at Lantian, in the warm sun, jade gives off smoke. \\
\zht{此情可待成追憶，} & This feeling --- need it wait to become a memory? \\
\zht{只是當時已惘然。} & Even then, at the time, I was already lost. \\
\bottomrule
\end{tabular}
\end{center}

\subsection{The allusions}

\begin{center}
\begin{tikzpicture}[c/.style={draw, thick, circle, fill=yellow!20, align=center, font=\small},
  al/.style={draw, rounded corners, align=center, text width=3.3cm, font=\scriptsize}]
  \node[c] (z) at (0,0) {\zht{錦瑟}\\ zither\\ 50 strings};
  \node[al, fill=blue!10] (a) at (-4.2,1.5) {\zht{莊周夢蝶} Zhuangzi dreams he is a butterfly: which is dreamer, which dream? \emph{illusion}};
  \node[al, fill=red!10] (b) at (4.2,1.5) {\zht{望帝杜鵑} a Shu king becomes the cuckoo whose cry is blood: \emph{grief, loss of a kingdom}};
  \node[al, fill=cyan!10] (c) at (-4.2,-1.5) {\zht{鮫人泣珠} mermen weep tears that turn to pearls; pearls wax with the moon: \emph{sorrow made beautiful}};
  \node[al, fill=green!10] (d) at (4.2,-1.5) {\zht{藍田玉煙} jade of Lantian, seen as smoke from afar, vanishes up close: \emph{the unreachable}};
  \draw[thick, gray] (z) -- (a); \draw[thick, gray] (z) -- (b);
  \draw[thick, gray] (z) -- (c); \draw[thick, gray] (z) -- (d);
  \node[font=\scriptsize, text=gray] at (0,2.3) {couplet 2};
  \node[font=\scriptsize, text=gray] at (0,-2.3) {couplet 3};
\end{tikzpicture}
\end{center}

\subsection{Interpretations}
\begin{itemize}
  \item \textbf{Love poem} for a lost woman (one tradition names a maid of Linghu Chu's
        household called Jinse).
  \item \textbf{Elegy for his wife}: the fifty strings broken into twenty-five (a legend
        says the Yellow Emperor had a fifty-string zither split in two because its music was
        too sad); a broken string (\zht{斷弦}) is a standard image for a wife's death.
  \item \textbf{Preface to his collected poems}, looking back on his whole life near fifty:
        each of the four middle images then names a mode of his poetry (dream, lament,
        sorrow, the elusive).
  \item \textbf{The key is the form}: the poem is about the impossibility of separating
        experience and memory; the meaning is ``lost'' (\zht{惘然}) even at the moment of
        feeling. The ambiguity is the subject.
\end{itemize}

\begin{remark}
Note how strictly regulated and perfectly parallel the middle couplets are (person /
person, dawn / spring, dream / heart, butterfly / cuckoo; sea / field, moon / sun, pearl /
jade, tears / smoke). Li Shangyin's obscurity is \emph{not} formal looseness: the clarity
of the structure makes the obscurity of the reference more haunting.
\end{remark}

\section{``Untitled'' \zht{無題} (``Hard it is to meet'')}

\begin{center}
\begin{tabular}{p{0.36\textwidth}p{0.54\textwidth}}
\toprule
\zht{相見時難別亦難，} & Hard it is to meet, and parting too is hard; \\
\zht{東風無力百花殘。} & the east wind has no strength, the hundred flowers wither. \\
\zht{春蠶到死絲方盡，} & The spring silkworm spins its thread until it dies; \\
\zht{蠟炬成灰淚始乾。} & the candle's tears dry only when it turns to ash. \\
\zht{曉鏡但愁雲鬢改，} & At dawn, in the mirror, she fears her cloud of hair has changed; \\
\zht{夜吟應覺月光寒。} & chanting poems at night, you must feel the moonlight's cold. \\
\zht{蓬山此去無多路，} & From here to Penglai Mountain is not so far a road: \\
\zht{青鳥殷勤為探看。} & blue bird, be diligent, go and look for me. \\
\bottomrule
\end{tabular}
\end{center}

\begin{itemize}
  \item \textbf{Puns}: \zht{絲} \emph{sī} ``silk'' = \zht{思} \emph{sī} ``longing''; candle
        wax drips as ``tears''. Couplet 2 became a proverb for devotion until death (now
        often applied to teachers!).
  \item \textbf{Shifting perspective}: line 5 seems to describe the woman, line 6 the man
        (or vice versa): the pronouns are absent.
  \item \textbf{Penglai} \zht{蓬山}: the isle of immortals; the \textbf{blue bird} is the
        messenger of the goddess Queen Mother of the West. The beloved is unreachable as an
        immortal (a Daoist nun?).
\end{itemize}

\section{``Night Rain: Sent North'' \zht{夜雨寄北}}

A rare plain poem by Li Shangyin, sent from Sichuan (Ba) to his wife or a friend in the
north.

\begin{center}
\begin{tabular}{p{0.24\textwidth}p{0.3\textwidth}p{0.36\textwidth}}
\toprule
\zht{君問歸期未有期，} & jūn wèn guī qī wèi yǒu qī & You ask when I'll come home: there's no date yet. \\
\zht{巴山夜雨漲秋池。} & Bā shān yè yǔ zhǎng qiū chí & Night rain in the hills of Ba floods the autumn pools. \\
\zht{何當共翦西窗燭，} & hé dāng gòng jiǎn xī chuāng zhú & When shall we trim the candle together at the west window \\
\zht{卻話巴山夜雨時。} & què huà Bā shān yè yǔ shí & and talk back over this night of rain in the hills of Ba? \\
\bottomrule
\end{tabular}
\end{center}

\begin{center}
\begin{tikzpicture}[t/.style={draw, thick, rounded corners, align=center, font=\scriptsize,
  text width=3.2cm, minimum height=1.1cm}]
  \node[t, fill=blue!10] (now) at (0,0) {\textbf{Present}\\ alone, rain in Ba\\ (lines 1--2)};
  \node[t, fill=orange!15] (fut) at (6,0) {\textbf{Imagined future}\\ together at the west window\\ (line 3)};
  \draw[-{Stealth}, thick] (now) to[bend left=25] node[above, font=\scriptsize] {``when shall we \ldots''} (fut);
  \draw[-{Stealth}, thick, dashed] (fut) to[bend left=25] node[below, font=\scriptsize] {``talk back over'' this night (line 4)} (now);
\end{tikzpicture}
\end{center}

\textbf{Analysis.} The phrase \zht{巴山夜雨} appears twice: first as present misery, then
as a future memory. The poem imagines the present from the future, so that loneliness
becomes something that will one day be shared. The repeated \zht{期} in line 1 is also
unusual in regulated verse, which normally avoids repetition: it sounds like speech.

\section*{Summary}
\begin{itemize}
  \item Late Tang: decline, nostalgia, sensuousness; Du Mu, Li Shangyin, Wen Tingyun.
  \item Du Mu's ``Mooring on the Qinhuai'': history in a pop song.
  \item Li Shangyin: caught between factions; ``Untitled'' poems; three reading strategies
        (biographical, allegorical, open).
  \item ``Ornamented Zither'': four allusions in perfectly parallel couplets; meaning lost
        ``even then''.
  \item ``Untitled'': silkworm and candle (silk = longing); ``Night Rain'': present seen
        from the future.
\end{itemize}
